% ---------------------------------------------------------------------
%  tfpt_figures.tex -- shared, reusable TikZ figure macros for every
%  TFPT document.  \input by tfpt_style.tex (which loads tikz + the
%  needed libraries + etoolbox).  Each macro draws a self-contained
%  centred figure; nothing is drawn until the macro is called.
%
%    \TFPTclaimstack[zone]  -- the architecture claim stack (B3); the
%                              optional zone in {axioms,anchor,compiler,
%                              sm,bridges,frontier} is highlighted.
%    \TFPTanchorcockpit      -- the anchor a=(1,1,2) cockpit (B4)
%    \TFPTeightglue          -- the D5 (+) A3 +mu4 => E8 glue (B5)
%    \TFPTmotif              -- the 2/3 motif map (B6)
%    \TFPTactionkfive        -- the K5 action-lattice mnemonic (B7)
% ---------------------------------------------------------------------

% ===== B3: the architecture claim stack =========================
\newcommand{\tfhl}[1]{\ifdefstring{\tfcur}{#1}{cbhi}{cbnl}}
\newcommand{\TFPTclaimstack}[1][none]{%
\def\tfcur{#1}%
\begin{center}
\begin{tikzpicture}[font=\footnotesize,
  cb/.style={rounded corners=2pt,inner sep=4pt,text width=118mm,align=center,minimum height=8.5mm},
  cbnl/.style={line width=0.5pt}, cbhi/.style={line width=1.5pt},
  ar/.style={-{Stealth[length=2.2mm]},tfgray,semithick}]
  \node[cb,fill=tfred!8,draw=tfred!75!black,\tfhl{axioms}] (axioms)
    {\textbf{P1}\,$c_3=\tfrac1{8\pi}$ \quad \textbf{P2}\,$g_{\mathrm{car}}=5$ \ \textcolor{tfgray}{\itshape(two inputs, [O] axioms)}};
  \node[cb,fill=tfblue!8,draw=tfblue,\tfhl{anchor},below=3.5mm of axioms] (anchor)
    {Anchor $a=(1,1,2)$:\ $e_1,e_2,e_3=4,5,2$ \ \textcolor{tfgray}{\itshape([E] exact)}};
  \node[cb,fill=tfgreen!10,draw=tfgreen!70!black,\tfhl{compiler},below=3.5mm of anchor] (compiler)
    {Discrete compiler:\ $D_5\oplus A_3+\mu_4\Rightarrow E_8$,\ $h=30$,\ $\varphi_0$ \ \textcolor{tfgray}{\itshape([E] exact)}};
  \node[cb,fill=tfgreen!7,draw=tfgreen!60!black,\tfhl{sm},below=3.5mm of compiler] (sm)
    {SM skeleton:\ three families, hypercharges, flavor matrix $R$ \ \textcolor{tfgray}{\itshape([E])}};
  \node[cb,fill=tfgold!12,draw=tfgold!80!black,\tfhl{bridges},below=3.5mm of sm] (bridges)
    {Physical bridges:\ $v_{\mathrm{geo}}$,\ $G_{\mathrm{net}}$,\ $F_{\mathrm{transfer}}$ \ \textcolor{tfgray}{\itshape([C] bridges)}};
  \node[cb,fill=tfgray!12,draw=tfgray!70!black,\tfhl{frontier},below=3.5mm of bridges] (frontier)
    {Frontier readouts:\ $\Lambda$, inflation, $\eta_B$, dark matter, QG \ \textcolor{tfgray}{\itshape([C]/[O])}};
  \draw[ar] (axioms)--(anchor); \draw[ar] (anchor)--(compiler);
  \draw[ar] (compiler)--(sm);   \draw[ar] (sm)--(bridges); \draw[ar] (bridges)--(frontier);
\end{tikzpicture}
\end{center}}

% ===== B4: the anchor cockpit ====================================
\newcommand{\TFPTanchorcockpit}{%
\begin{center}
\begin{tikzpicture}[font=\footnotesize,
  pan/.style={rounded corners=2pt,draw=tfblue!60,fill=tfblue!4,inner sep=5pt,
    text width=46mm,align=left},
  hd/.style={font=\bfseries\small,tfblue}]
  \node[pan] (e) {\textbf{elementary} $e_k(a)$\\[2pt]
     $e_1=4\to|\mu_4|$\\ $e_2=5\to g_{\mathrm{car}}$\\ $e_3=2\to|\mathbb Z_2|$\\[2pt]
     $c_3=\tfrac1{2e_1\pi}=\tfrac1{8\pi}$};
  \node[pan,right=4mm of e] (p) {\textbf{power sums} $p_n=2+2^n$\\[2pt]
     $p_0,p_1,p_2,p_3=3,4,6,10$\\ $p_1p_2p_3=240=|R(E_8)|$\\ $p_4-p_3=8=\operatorname{rank}E_8$\\
     $\Rightarrow\dim E_8=248$};
  \node[pan,right=4mm of p] (m) {\textbf{derived motifs}\\[2pt]
     $4,5,8,16,25,$\\ $30,41,120,240$\\[2pt]
     \textcolor{tfgray}{\itshape one anchor,}\\ \textcolor{tfgray}{\itshape many projections}};
  \node[hd,above=1mm of p.north] {$a=(1,1,2)$ \ \textcolor{tfgray}{\footnotesize\itshape[E]}};
\end{tikzpicture}
\end{center}}

% ===== B5: the E8 glue ===========================================
\newcommand{\TFPTeightglue}{%
\begin{center}
\begin{tikzpicture}[font=\footnotesize,
  d/.style={rounded corners=2pt,draw=tfblue,fill=tfblue!6,inner sep=4pt,align=center,text width=36mm},
  mid/.style={rounded corners=2pt,draw=tfgold!80!black,fill=tfgold!12,inner sep=4pt,align=center,text width=52mm},
  e8/.style={rounded corners=2pt,draw=tfgreen!70!black,fill=tfgreen!10,inner sep=4pt,align=center,text width=52mm,font=\bfseries},
  rt/.style={rounded corners=2pt,draw=tfgreen!55!black,fill=tfgreen!6,inner sep=3pt,align=center,text width=52mm},
  side/.style={rounded corners=2pt,draw=tfgray!55,fill=tfgray!6,inner sep=4pt,align=left,text width=34mm,font=\scriptsize},
  ar/.style={-{Stealth[length=2mm]},tfgray,semithick}]
  \node[d] (d5) {$D_5$ discriminant\\ $\mathbb Z_4$, $q=5/4$};
  \node[d,right=10mm of d5] (a3) {$A_3$ discriminant\\ $\mathbb Z_4$, $q=3/4$};
  \node[mid,below=5mm of $(d5)!0.5!(a3)$] (glue) {$\mu_4$ anti-isometry\ ($q$-sum $=2$)};
  \node[e8,below=4mm of glue] (e8) {even unimodular $E_8$ lattice};
  \node[rt,below=4mm of e8] (rt) {$240$ roots,\ rank $8$,\ $\dim=248$};
  \node[side,right=8mm of glue] (why) {\textbf{why unique?}\\ rank $8$\\ even\\ unimodular\\ positive definite};
  \draw[ar] (d5)--(glue); \draw[ar] (a3)--(glue);
  \draw[ar] (glue)--(e8); \draw[ar] (e8)--(rt);
\end{tikzpicture}\\[2pt]
{\scriptsize\itshape $D_5\oplus A_3+\mu_4\Rightarrow E_8$ --- the unique even unimodular rank-$8$ glue \ [E] exact (lattice).}
\end{center}}

% ===== B5b: the double pushout (one glue, two categories) ========
\newcommand{\TFPTdoublepushout}{%
\begin{center}
\begin{tikzpicture}[font=\footnotesize,>={Stealth[length=2mm]},
  lat/.style={rounded corners=2pt,draw=tfblue,fill=tfblue!6,inner sep=5pt,align=center,text width=34mm},
  net/.style={rounded corners=2pt,draw=tfblue!70!black,fill=tfblue!4,inner sep=5pt,align=center,text width=34mm},
  e8l/.style={rounded corners=2pt,draw=tfgold!80!black,fill=tfgold!12,inner sep=5pt,align=center,text width=34mm,font=\bfseries},
  e8n/.style={rounded corners=2pt,draw=tfgreen!70!black,fill=tfgreen!10,inner sep=5pt,align=center,text width=34mm,font=\bfseries},
  ar/.style={->,tfgray,semithick},
  lb/.style={font=\scriptsize,text=black,align=center,fill=white,inner sep=1.5pt}]
  \node[lat] (LL) at (0,2.5)   {$D_5\oplus A_3$\\[1pt]\scriptsize root lattices,\ $\mathrm{disc}=\Z_4$};
  \node[e8l] (LR) at (6.6,2.5) {$E_8$\\[1pt]\scriptsize even unimodular,\ $\det 1$};
  \node[net] (BL) at (0,0)     {$(D_5)_1\otimes(A_3)_1$\\[1pt]\scriptsize chiral nets,\ $c=8$};
  \node[e8n] (BR) at (6.6,0)   {$(E_8)_1$\\[1pt]\scriptsize holomorphic,\ $\mu$-index $1$};
  \draw[ar] (LL)--node[lb]{$+\,\mu_4$ glue}(LR);
  \draw[ar] (BL)--node[lb]{$\rtimes\,\mu_4$ extension}(BR);
  \draw[ar] (LL)--node[lb,left=-1pt]{level-$1$}(BL);
  \draw[ar] (LR)--node[lb,right=-1pt]{level-$1$}(BR);
\end{tikzpicture}\\[2pt]
{\scriptsize\itshape One operation, two categories: the order-$|\mu_4|{=}4$ glue is the lattice pushout
$D_5\oplus A_3+\mu_4\Rightarrow E_8$ \emph{and} the simple-current / anyon-condensation extension
$(D_5)_1\otimes(A_3)_1\rtimes\mu_4\cong(E_8)_1$; the square commutes. \ [E] exact.}
\end{center}}

% ===== B6: the 2/3 motif map =====================================
\newcommand{\TFPTmotif}{%
\begin{center}
\begin{tikzpicture}[font=\footnotesize,
  ex/.style={rounded corners=2pt,draw=tfgreen!70!black,fill=tfgreen!8,inner sep=4pt,align=center,text width=92mm},
  br/.style={rounded corners=2pt,draw=tfgold!80!black,fill=tfgold!12,inner sep=4pt,align=center,text width=92mm},
  ar/.style={-{Stealth[length=2mm]},tfgray,semithick}]
  \node[ex] (a) {$|\mathbb Z_2|/N_{\mathrm{fam}}=2/3$ \ \textcolor{tfgray}{\itshape(anchor ratio)}};
  \node[ex,below=3mm of a] (b) {parabolic weights $\{0,\tfrac13,\tfrac23\}=\operatorname{spec}A_0^\star$};
  \node[ex,below=3mm of b] (c) {transfer spectrum $\{1,(2/3)^6,(1/3)^6\}$,\ gap $6\log\tfrac32$};
  \node[br,below=3mm of c] (d) {Koide target $2/3$ \ \textcolor{tfgray}{\itshape(source$\to$pole bridge [C])}};
  \node[br,below=3mm of d] (e) {Nariai / horizon entropy bound $2/3$ \ \textcolor{tfgray}{\itshape([C])}};
  \draw[ar] (a)--(b); \draw[ar] (b)--(c); \draw[ar] (c)--(d); \draw[ar] (d)--(e);
\end{tikzpicture}\\[2pt]
{\scriptsize\itshape Green: exact motif ([E]).\quad Gold: physical bridge ([C]) --- the same ratio, two epistemic types.}
\end{center}}

% ===== B7: the K5 action-lattice mnemonic ========================
%  Layout: complete graph K5 (pentagon, all 10 edges) on the LEFT, the
%  Pascal-row reading list on the RIGHT, vertically centred on the graph
%  so the two blocks never overlap.
\newcommand{\TFPTactionkfive}{%
\begin{center}
\begin{tikzpicture}[font=\scriptsize]
  % --- K5 vertices on a regular pentagon, centred at the origin ---
  \foreach \i in {1,...,5}{\coordinate (k\i) at ({90+72*(\i-1)}:12mm);}
  \begin{scope}[on background layer]
    \foreach \i in {1,...,5}{\foreach \j in {1,...,5}{\ifnum\i<\j
      \draw[tfgray!55,line width=0.5pt] (k\i)--(k\j);\fi}}
  \end{scope}
  \foreach \i in {1,...,5}{\fill[tfblue!70!black] (k\i) circle (1.8pt);}
  \node[tfgray,font=\scriptsize\itshape] at (0,-1.9) {$K_5$};
  % --- reading list, anchored to the right of the whole graph, centred ---
  \node[anchor=west,align=left,inner sep=0pt] at (1.9cm,0) {%
    $\tbinom51=5$\ \ vertices --- action $x^1$ (Hubble slot)\\[1pt]
    $\tbinom52=10$\ edges --- $\Lambda$ density pair $x^{10}$\\[1pt]
    $\tbinom53=10$\ triangles --- dual sector\\[1pt]
    $\tbinom54=5$\ \ complements --- inverse sector\\[1pt]
    $\tbinom55=1$\ \ determinant --- closure};
\end{tikzpicture}\\[2pt]
{\scriptsize\itshape $K_5$ on the five carrier slots; the action powers $x^{1},x^{5},x^{10}$ are its
$\tbinom5k$ faces \ \textcolor{tfgray}{(mnemonic; the physical scale map stays [C]).}}
\end{center}}
